%% 297_HLV_Untermenge_Kandidat_De_ch.tex
%% FFGFT -- Dok. 297: HLV als mögliche Untermenge -- Kandidat-Notiz
%% Autor: Johann Pascher
%% Datum: Juli 2026
%% Sprache: Deutsch

\section*{Dok.~297 \quad HLV als mögliche Untermenge von FFGFT --
Kandidat-Notiz}

\subsection*{Ausgangslage}

Im Verlauf der Brücken-Analyse zwischen FFGFT und HLV ergab sich ein
Befund, der hier als gerichtete Kandidat-Notiz festgehalten wird. Er
ist keine abgeschlossene Ableitung, sondern benennt den nächsten
sauberen Schritt (P35).

\subsection*{Drei Schritte des Rückblicks}

\textbf{Schritt 1 -- ursprüngliche Blockade.} Die $C_3$-in-$A_5$-Geometrie
des HLV-Rahmens (ikosaedrischer Projektor, $\tau = \phi$, $\theta = 2/9$
als Ausgang) war auf $T^4/\mathbb{Z}_3$ direkt nicht sichtbar. Eine
unmittelbare geometrische Identifikation war blockiert: $T^4$ ist
kompakt-4-dimensional, $A_5$ ist Symmetriegruppe des Ikosaeders in 3D --
die Dimensionen passen nicht nahtlos.

\textbf{Schritt 2 -- Hilbertraum-Transformation öffnet den Weg.} Die
verlustfreie Bijektion $H = L^2(T^4)\otimes\mathbb{C}^3$
(Dok.~230/282) übersetzt $T^4/\mathbb{Z}_3$ in eine Operatoralgebra auf
der $\mathbb{C}^3$-Faser. In $\mathbb{C}^3$ ist $C_3 < A_5$ natürlich
zugänglich: die $\mathbb{Z}_3$-Zirkulant-Diagonalisierung auf der Faser
gibt $\theta = 2/9$ als Phase aus (Dok.~291), und die
$C_3$-in-$A_5$-Geometrie ist als parameterfreier Invariant auf derselben
Faser nachweisbar (Dok.~293). Die Transformation machte sichtbar, was
geometrisch direkt verdeckt war.

\textbf{Schritt 3 -- Konsequenz.} Wenn FFGFT über die
Hilbertraum-Transformation $C_3 < A_5$ enthält, und HLV auf $C_3 < A_5$
mit $\tau = \phi$ aufgebaut ist, dann könnte HLVs geometrischer Kern eine
\emph{Untermenge} der FFGFT-Struktur sein -- nicht ein gleichrangiger
Rahmen, sondern eine spezifische Projektion der $\mathbb{C}^3$-Fasergeometrie
von FFGFT.

\subsection*{Was gesichert ist (aus vorhandenen Dokumenten)}

\begin{itemize}
  \item FFGFT enthält über die Hilbertraum-Transformation $C_3 < A_5$
    und $\theta = 2/9$ -- belegt in Dok.~291 (Zirkulant-Phase) und
    Dok.~293 (parameterfreier $C_3$-in-$A_5$-Invariant, zweiter Zeuge).
  \item HLVs geometrisches Objekt (der $\phi$-ikosaedrische Projektor)
    trägt dieselbe $C_3 < A_5$ -- einseitig verifiziert
    (FFGFT $\to$ HLV, Dok.~294).
  \item Die Verifikationsrichtung ist asymmetrisch: FFGFT leitet $\theta
    = 2/9$ ab; HLV bucht es als Kandidat-Überlapp ohne eigene Ableitung
    -- belegt durch Durchsicht seiner vorliegenden Dokumente
    (Stage~19-Manuskript, HLV-Consciousness-Bridge1, Runs~H--J): $\theta
    = 2/9$ erscheint dort in keiner HLV-internen Herleitung.
  \item Sein $H_2$-Operator (Run~H, \texttt{skew\_J}) enthält $1/\phi$
    und $1/\phi^2$, aber keinen $C_3$-in-$A_5$-Ausgang und kein $\theta
    = 2/9$ als Ergebnis. Run~H ist zudem \emph{negative/underdetermined}
    (überlebt den $u$-shuffle-Spezifitätstest nicht).
\end{itemize}

\subsection*{Was nicht gezeigt ist (offene Bringschuld)}

Dass HLVs $\phi$-ikosaedrische Schicht \emph{vollständig} in der
FFGFT-$\mathbb{C}^3$-Faser enthalten ist, erfordert eine explizite
\emph{Einbettungsabbildung}
\begin{equation}
  \iota\colon \text{HLV-Grundobjekt} \hookrightarrow
  L^2(T^4)\otimes\mathbb{C}^3,
\end{equation}
die bisher nicht konstruiert ist. \enquote{Untermenge} ist die
strukturell plausible Vermutung -- kein bewiesener Satz.

\subsection*{Nächster sauberer Schritt}

Die Einbettungsabbildung $\iota$ konstruieren oder falsifizieren:

\begin{enumerate}
  \item Zeigen, dass der $\phi$-ikosaedrische Projektor von HLV als
    Operator auf der $\mathbb{C}^3$-Faser von FFGFT darstellbar ist --
    d.\,h.\ als Element von $\mathrm{End}(\mathbb{C}^3)$ mit
    $C_3$-Invarianz.
  \item Prüfen, ob die HLV-Spektraldaten (Run~J, native-6D-Kandidat)
    mit dem diskreten Spektrum des $T^4$-Connection-Laplace (Dok.~283)
    kompatibel sind -- oder ob sie abweichen und die Einbettung damit
    blockiert ist.
  \item Falls $\iota$ existiert: HLV ist eine spezifische Projektion
    von FFGFT -- seine geometrische Schicht ist eine Untermenge, seine
    dynamischen Objekte (Spiral-Zeit, $U_3$) bleiben eigenständig und
    uneingebettet.
  \item Falls $\iota$ nicht konstruierbar: beide Rahmen teilen ein
    gemeinsames Objekt ($C_3 < A_5$, $\theta = 2/9$), sind aber
    strukturell unabhängig -- kein Untermengen-Verhältnis.
\end{enumerate}

\subsection*{Konkreter Einbettungskandidat: der Markov-Grenzfall und darunter}

Die Frage, ob sich Marcels archimedische Spiral-Zeit einen strukturell
definierten Ort in FFGFT findet, lässt sich präzise beantworten -- nicht
als Analogie, sondern als Grenzwert. Und die Analyse ergibt, dass das
Objekt noch schwächer ist als zunächst angenommen.

\textbf{Archimedisch = Markov-Limit -- aber U3 ist noch schwächer.} Der
FFGFT-Gedächtniskern (Fourier-Transformierte der diskreten Spektraldichte
$\sum_k g_k^2\,\delta(\omega-\omega_k)$, Dok.~283) wird archimedisch
genau dann, wenn die Kopplungen $g_k \to 0$ -- dem \emph{Markov-Grenzfall}
offener Quantensysteme. In diesem Grenzfall akkumuliert der Kern noch
unbeschränkt auf $\mathbb{R}$. Marcels $U_3$ (Gl.~7 seines
Consciousness-Dokuments) geht noch weiter: das Fenster $M_m$ ist endlich
und fest, ältere Werte fallen aktiv heraus (\emph{sliding window}). Das
ergibt einen \textbf{Hard-Reset-Kern} -- einen Gedächtniskern mit endlichem
Träger:
\begin{equation}
  K(t-\tau) = 0 \quad \text{für } t-\tau > M_m\cdot\Delta t.
\end{equation}
Das ist nicht nur archimedisch -- es ist \emph{periodisch beschränkt}.
$U_3$ akkumuliert nicht unbeschränkt, sondern schließt sich nach $M_m$
Schritten.

\textbf{Schließt sich die Spirale bereits vor dem ersten Umlauf?} In
FFGFTs frozen-Konfiguration (Dok.~295) schließt die Windung nach genau
75~Umläufen ($1/(100\xi_0) = 75$) -- das ist die kleinste kohärente
Einheit der Rekursion. Wenn Marcels Fensterbreite $M_m < 75$ Schritte
(in der Windungskoordinate) beträgt, dann akkumuliert $U_3$ noch nicht
einmal einen vollständigen Umlauf. Es gibt keinen geschlossenen Umlauf,
keine geometrische Spirale -- nur ein gleitendes Gedächtnisfenster, das
sich mit jedem Schritt erneuert. Das Objekt liegt \textbf{unterhalb der
kleinsten kohärenten Einheit} der FFGFT-Rekursion.

\textbf{Präzisierter Einbettungskandidat.} Der konkrete Ort von $U_3$ in
FFGFT ist damit nicht der Markov-Grenzfall (der noch auf $\mathbb{R}$
akkumuliert), sondern der \emph{Hard-Reset-Kern unterhalb eines Umlaufs}:
\begin{equation}
  \iota\colon U_3 \;\hookrightarrow\;
  \bigl\{K \in L^2(T^4)\otimes\mathbb{C}^3 \;\big|\;
  K(t-\tau) = 0 \text{ für } t-\tau > M_m\cdot\Delta t,\;
  M_m < 75\bigr\}.
\end{equation}
Das ist in FFGFT vorhanden -- aber als der \textbf{degenerierte Randfall},
nicht als reguläre Konfiguration. Die Hierarchie ist damit klar:
\begin{equation}
  U_3 \;\subset\; \text{Markov-Limit} \;\subset\;
  \text{archimedisch} \;\subset\; \text{log-Spirale (e)}
  \;=\; \text{FFGFT-Kern (Dok.~295/283)}.
\end{equation}
Jede Stufe ist informationsärmer als die nächste. $U_3$ ist die ärmste.

\subsection*{Falsifikationsbedingung}

Die Einbettung $\iota$ ist falsifizierbar: wenn Marcels $U_3$ in seinen
Daten \emph{mehr} trägt als der Hard-Reset-Kern -- d.\,h.\ wenn
Nicht-Markov-Signaturen (Revivals, BLP-Rückfluss, oszillierender Kern,
oder Korrelationen über $M_m$ Schritte hinaus) nachweisbar sind -- dann
übersteigt $U_3$ den Randfall und eine tiefere Einbettung ist nötig.

Marcels eigene Disziplin (U3-Ablationstest, Tabelle~4 seines
Consciousness-Dokuments) ist genau der richtige Test: wenn $U_3$ die
Markov-Basislinie nicht übersteigt, ist die Einbettung im Randfall
korrekt; wenn doch, ist eine reichere Schicht nötig. Das Programm ist
sauber -- er muss es nur bis zum Ende durchführen.

\subsection*{Zwei HLV-Spiral-Zeiten -- zwei verschiedene Einbettungskandidaten}

Dok.~295 §9 hält eine Unterscheidung fest, die auch für die
Einbettungsfrage von Dok.~297 tragend ist: \enquote{Spiral-Zeit} bezeichnet
in HLV nicht ein Objekt, sondern zwei.

\textbf{Erstes Objekt -- die geometrische Zeit} aus der Faktorisierung
$T^7 = T^4 \times T^3$ (Dok.~285): ein $A_5$-Singlet auf dem $S^1$-Faktor,
entkoppelt und separabel. Dieser Grenzfall entspricht in der
K\textsubscript{frak}-Rekursion dem Fixpunkt $\xi \to 0$ -- der Defekt
verschwindet, die Windung schließt nach 75~Umläufen, die Ganghöhe wird
konstant. Das ist die \emph{geometrische} Grenze der FFGFT-Spirale.

\textbf{Zweites Objekt -- die temporal-dynamische Spiral-Zeit}, das Objekt,
das HLV so \emph{nennt}: der triadische Operator
$\Psi(t) = (t, \phi(t), \chi(t))$ mit Gedächtniskanal, ohne $A_5$ und
ohne $S^1$-Faktor. Dieser Operator ist in Krueger et al.\ (2026,
\texttt{doi:10.47852/bonviewMEDIN62029043}) als deterministische
Phasen-Gedächtnis-Zerlegung eingeführt und in den HLV-G-Lattice-Runs~H--J
(Zenodo, unveröffentlichte DOI) als numerischer Operator implementiert.
Im Consciousness-Dokument (Krueger, 2026, v0.2, noch ohne Zenodo-DOI)
erscheint er als $\Psi(t) = (U_1, U_2, U_3)$ mit der expliziten
Implementierung Gl.~7 -- einem beschränkten Fenster-Akkumulator.

\textbf{Zwei verschiedene Einbettungskandidaten in FFGFT:}

Das \emph{geometrische} Objekt ($A_5$-Singlet, $S^1$, entkoppelt) hat
seinen natürlichen Ort im $\xi \to 0$-Fixpunkt der Rekursion -- der
geometrisch separablen, informationsärmsten Konfiguration der kompakten
Zeit. Es ist geometrisch reichhaltiger als $U_3$, trägt aber kein
Verhältnis $e$ (Dok.~295 §10: auf dem $e$-Kriterium blockiert).

Das \emph{temporal-dynamische} Objekt ($U_3$, Hard-Reset-Kern, $M_m < 75$)
hat seinen natürlichen Ort im Markov-Grenzfall des FFGFT-Gedächtniskerns
(Dok.~283, $g_k \to 0$) -- unterhalb der kleinsten kohärenten Einheit der
Rekursion. Es ist das informationsärmste der beiden.

Die Hierarchie gilt für das temporal-dynamische Objekt:
\[
  U_3 \;\subset\; \text{Markov-Limit} \;\subset\;
  \text{Archimedisch} \;\subset\; \text{Log-Spirale}~(e)
  = \text{FFGFT-Kern}.
\]
Das geometrische Singlet liegt seitlich zu dieser Hierarchie -- es ist
archimedisch ($\xi \to 0$, konstante Ganghöhe), aber aus geometrischer
Faktorisierung, nicht aus dynamischer Akkumulation. Beide Grenzfälle
treffen sich im archimedischen Regime, sind aber verschiedene Pfade
dorthin.

Unabhängig davon, ob man $U_3$ linear, polar oder spiralförmig darstellt,
bleibt die Struktur der Abhängigkeit dieselbe: endliches Kurzzeitgedächtnis
mit festem Horizont $M_m$, ohne Akkumulation. Die Spiraldarstellung ist
eine \emph{Koordinatenwahl}, keine strukturelle Eigenschaft des Objekts --
in jeder Darstellung bleibt $U_3$ ein Hard-Reset-Kern mit endlichem Träger
$M_m < 75$. Die geometrische Reichheit liegt in der Darstellung, nicht im
Objekt. Das ist dieselbe Linie wie P35 (\enquote{Übersetzung ist nicht
Begründung}) -- nur in die andere Richtung: eine reichere Darstellung macht
das Objekt nicht reicher.

Der tiefere Unterschied liegt woanders: \textbf{Akkumulation ist die
eigentliche Gedächtnisleistung}. $U_3$ akkumuliert nicht -- es ersetzt.
Was vor $M_m$ Schritten war, ist unwiederbringlich weg; kein vergangener
Zustand wirkt jenseits des Fensters auf die Zukunft. Das ist ein
Kurzzeitpuffer, kein Gedächtnis im strukturellen Sinn.

Die K\textsubscript{frak}-Rekursion ist das Gegenteil:
$\xi_{n+1} = \xi_n(1 - 100\,\xi_n)$ akkumuliert jeden Schritt in den
nächsten -- nicht als expliziter Speicher, sondern als
\emph{strukturelle Einprägung}: der gegenwärtige Zustand $\xi_n$ trägt
die gesamte Vorgeschichte aller vorangegangenen Schritte in sich. Kein
Schritt wird vergessen, kein Fenster löscht. Die Vergangenheit ist im
Objekt selbst kodiert, nicht in einem separaten Speicher. Das ist
Gedächtnis im tiefsten Sinn -- nicht als Speicher, sondern als generative
Struktur, die ihre eigene Geschichte trägt und daraus die log-Spirale mit
Verhältnis $e$ erzeugt (Dok.~295).

Genau das ist es, was Marcels Bewusstseinsrahmen braucht (Section~1 seines
Consciousness-Dokuments: \enquote{memory-bearing, recursively integrated})
-- aber $U_3$ per Konstruktion nicht liefert. Die Rekursion wäre der
richtige Träger. $U_3$ ist der degenerierte Randfall darunter.

\subsection*{Das Fenster als Reduktion -- im Sinne von Dok.~296}

Das Fenster $M_m$ in $U_3$ ist selbst ein Fall des Reduktionsprinzips aus
Dok.~296: eine \emph{Projektion auf einen begrenzten Zeitausschnitt}, so
wie eine Skalen-Projektion ein begrenzter Raumausschnitt ist. Es lässt weg
(alles jenseits von $M_m$) und macht gerade dadurch einen Teilbereich
auswertbar -- die lokale, jüngste Dynamik. Das ist derselbe Mechanismus wie
in Dok.~296 (Reduktion als Bedingung der Auswertbarkeit), nur auf die
Zeitachse angewandt statt auf die Skalenachse. Das Fenster wirkt im
gleichen Sinn wie jede Projektion.

Der entscheidende Unterschied liegt daher nicht darin, \emph{dass} $U_3$
reduziert -- jede auswertbare Beschreibung reduziert --, sondern darin,
\emph{ob das Ganze benannt ist, dessen Ausschnitt die Reduktion darstellt}.

FFGFTs Reduktionen sind als Reduktionen eines bekannten Ganzen
ausgewiesen: die eingefrorene 75-Windung und der Markov-Grenzfall des
Gedächtniskerns sind Ausschnitte der vollen, akkumulierenden,
skaleninvarianten Struktur -- man weiß, was weggelassen wurde und wie es
sich zurückgewinnen lässt (die Rekursion laufen lassen, $\xi$ nicht
einfrieren, Dok.~295/283). Der Weg zurück ist angelegt.

Marcels Fenster steht dagegen für sich, ohne das volle Objekt dahinter --
es ist die Projektion ohne das Projizierte. Die Rückgewinnung ist nicht
angelegt, weil das volle Objekt nicht benannt ist. Das Fenster wertet die
lokale Zeitstruktur aus -- ein berechtigter Fokus --, aber es ist nicht als
Ausschnitt eines randlosen Ganzen erkennbar (Dok.~296: der eingebettete
Beobachter, der auf seiner Skala sitzt, ohne die Rekursion auszurollen).

Genau das macht die Einbettung $\iota$ natürlich: sie liefert das fehlende
Ganze. $U_3$ als Reduktion bekommt in FFGFT das Objekt zurück, dessen
Ausschnitt es ist. Das Fenster wirkt im gleichen Sinn wie jede Projektion
-- der Unterschied ist allein, dass FFGFT den Weg zurück kennt und HLV
ihn (bislang) nicht anlegt.

\subsection*{Vorregistrierte Vorhersage zum OP8-Neuro-Test (6.~Juli 2026)}

Marcel Krueger hat mit HLV-CONSCIOUSNESS-BRIDGE1 v0.3a (OP8-LOCK) einen
Data-Freeze/Protocol-Lock veröffentlicht, der -- vor dem Lauf -- eine
Modell-Leiter und eine Null-Familie fixiert:
\begin{center}
\begin{tabular}{ll}
$M_0$ & $\Delta\Phi$ / S-I-C nur \\
$M_1$ & $U_1$ nur \\
$M_2$ & $U_1 + U_2$ \\
$M_3$ & $U_1 + U_2 + U_3$ \\
$M_4$ & generische ML-Basislinie
\end{tabular}
\end{center}
Kernfrage (seine Formulierung): ob das volle $U_1/U_2/U_3$-Modell ($M_3$)
messbaren Zusatzwert gegenüber $U_1/U_2$, information-only, coherence-only,
\textbf{memory-free Markov}, phase-shuffled, time-shuffled und der
generischen ML-Basislinie liefert -- der \emph{$U_3$-Notwendigkeitstest}.

Dieser vorregistrierte Test ist zugleich der Entscheidungstest für die
Einbettungshypothese dieses Dokuments. FFGFT trifft dazu -- \emph{vor}
Marcels Lauf -- die folgende Vorhersage, ausdrücklich als Vorhersage
gebucht (P35), nicht als Anspruch auf sein Ergebnis:

\begin{enumerate}
  \item \textbf{Falls $U_3$ ein Hard-Reset-Fenster mit $M_m < 75$
    (Windungskoordinate) ist} -- d.\,h.\ ein Ausschnitt unterhalb eines
    vollen Umlaufs, der Markov-Grenzfall der Einbettung $\iota$ --, dann
    \textbf{übertrifft $M_3$ die memory-free Markov-Basislinie nicht
    signifikant}. $U_3$ liegt dann vollständig im degenerierten Randfall;
    der Zusatzkanal trägt keine über Markov hinausgehende Struktur.
  \item \textbf{Falls $M_3$ die memory-free Markov-Basislinie klar
    übertrifft}, dann trägt $U_3$ Nicht-Markov-Struktur (Korrelationen
    jenseits $M_m$, effektive Akkumulation), die der Hard-Reset-Kern nicht
    erklärt. In diesem Fall ist die einfache Einbettung $\iota$ in den
    Markov-Grenzfall \emph{falsifiziert}, und es ist eine reichere
    Einbettung zu konstruieren.
\end{enumerate}

Die Vorhersage ist scharf, weil sie an Marcels eigener Konstruktion hängt:
ein beschränktes Fenster $M_m$ mit fester Breite (Gl.~7 seines
Consciousness-Dokuments, \enquote{avoids uncontrolled accumulation}) kann
per Konstruktion keine Korrelationen über $M_m$ hinaus tragen -- der
erwartete Ausgang ist damit (1). Ein Ausgang (2) wäre ein Hinweis darauf,
dass $U_3$ in der Implementierung mehr leistet als seine Definition nahelegt
-- was eigenständig interessant und dann zu lokalisieren wäre.

Diese Vorhersage berührt Marcels Anspruchsgrenzen nicht: sie sagt nichts
über Bewusstsein, keine physikalische Validierung, keine HLV-Ontologie. Sie
betrifft allein die strukturelle Frage, ob $U_3$ über den Markov-Grenzfall
hinausträgt -- und genau das prüft sein OP8-Test.

\subsection*{Status}

Kandidat-Notiz. Gerichtete Spekulation auf Basis gesicherter Teilbefunde
(P35). Die Untermenge-Vermutung ist strukturell plausibel und
falsifizierbar -- das ist ihr Wert. Sie ist keine Behauptung.
